\begin{helper}
The unchanged summand comparison across one split-blocker cut is obtained by
finite additivity over the candidate-indexed union, not by comparing one
unspecified candidate cylinder.  In any fixed actual witness surface, suppose
the unchanged source event and unchanged target event are written as finite
unions
\[
U^-=\bigcup_{x\in B_k} U^-_x,\qquad
U^+=\bigcup_{x\in B_k} U^+_x,
\]
where the candidate pieces are pairwise almost-disjoint on each side,
measurable up to null sets, and have equal source and target masses:
\(\nu(U^-_x)=\nu'(U^+_x)\) for every candidate \(x\in B_k\).  Then
\(\nu(U^-)=\nu'(U^+)\).
\end{helper}
\begin{proof}
Replace the two unchanged events by their displayed finite candidate unions.
Because the source pieces are pairwise almost-disjoint and null-measurable,
finite additivity gives
\[
\nu(U^-)=\sum_{x\in B_k}\nu(U^-_x).
\]
The same argument on the target side gives
\[
\nu'(U^+)=\sum_{x\in B_k}\nu'(U^+_x).
\]
The summands agree candidate by candidate by hypothesis.  Therefore the two
finite sums agree, and hence the two unchanged-event masses agree.  This is
the exact additivity step needed after the finite product/cylinder and
marginal calculations have supplied equality on each candidate cell; overlaps
are handled by the stated almost-disjointness hypotheses rather than ignored.
\end{proof}
